% TOP_PROBLEM: {"id": 30004281, "title": "Small-Weight Quantum Operations on Lattice Conformal Nets", "category_id": 16, "category": {"id": 16, "name": "physics", "display_name": "Mathematical Physics", "description": "Problems at the intersection of mathematics and physics.", "slug": "physics", "order_index": 16, "created_at": "2026-07-31T15:26:25.671Z"}, "status": "open", "statusQualification": "Primary Problem 9 (PDF page 30, printed page 3082) asks for small positive real weights $w>1$. The catalogue adds arbitrarily small, which can be read as weights approaching $1$; that stronger quantifier was not verified in the primary source. The original small-weight construction question is preserved as source context, and the inherited stronger wording is explicitly not endorsed as a recovered theorem target. The nontriviality and actual weight convention exclude the vacuous identity operation."}
% FIRST_POSTED: 2026-10-11
% ENTRY_KIND: statement_only
% SUBMISSION: {"kind": "statement_only", "category": "top_problems", "problem_id": "30004281", "proposal_type": "established", "source": "Subfactors and Applications (2019), Oberwolfach Reports, DOI 10.4171/owr/2019/49, pp. 3053--3140. Primary report accessed 2026-10-11; relevant PDF pages 29, 30. https://doi.org/10.4171/owr/2019/49; Bischoff, generalized orbifold construction for conformal nets (catalogue-listed context; not independently consulted in this preparation). https://arxiv.org/abs/1608.00253; Osborne et al., numerical evidence for a Haagerup CFT (catalogue-listed context; not independently consulted in this preparation). https://authors.library.caltech.edu/records/mgvpx-37124", "status_evidence": "Canonical source status open retained. Supplied literature triage dated 2026-08-22: The hypergroup and generalized-orbifold framework supplies examples and structural constraints, and numerical work gives evidence for Haagerup-related CFT behavior, but the requested quantum operation for the lattice or VOA examples, especially the Haagerup case, has not been constructed. Primary-report access checked 2026-10-11; this historical report and inherited triage do not certify current resolution of every component.", "status_checked_on": "2026-10-11", "duplicate_check": "Compared canonical ID, title and normalized target with the complete local UnsolvedMath registry; normalized exact-text matches: none; title matches: none. Related or overlapping canonical records: none identified in this batch. Primary source context and the existing OWR statement audit used for scope and duplicate review; an exhaustive mathematical-equivalence/literature search was not performed.", "attribution": "Original proposer not named in the supplied canonical record; attribution remains to the identified Oberwolfach contribution/report, not its editors as invented proposers. Explanatory statement preparation and source/semantic audit assisted by OpenAI GPT-6 Codex; no human mathematical authorship or independent proof verification is claimed."}
% SOURCE_STATEMENT: "For a lattice conformal net or a strongly local vertex-operator-algebra net $\\mathcal A$, find $\\phi\\in\\operatorname{QuOp}(\\mathcal A)$ satisfying $$\\phi\\circ\\phi=w^{-1}\\operatorname{id}+(1-w^{-1})\\phi$$ for arbitrarily small real $w>1$."
% Preparation and final self-audit: CONTRIBUTING.md prompts C and D; source/audit artifacts retained in the local physics_statements_2026_10_11 batch.
% Imported from: physics_statements.tex; UnsolvedMath ID 30004281
% !TeX program = lualatex
\documentclass[11pt]{article}
\usepackage[margin=25mm]{geometry}
\usepackage{fontspec}
\setmainfont{Latin Modern Roman}
\usepackage{amsmath,amssymb,mathtools}
\usepackage{unicode-math}
\setmathfont{Latin Modern Math}
\usepackage{xurl,hyperref}
\hypersetup{colorlinks=true,urlcolor=blue}
\setcounter{secnumdepth}{0}
\setlength{\parindent}{0pt}
\setlength{\parskip}{5pt}
\setlength{\emergencystretch}{3em}
\newcommand{\sref}[2]{\hyperlink{source:#1:#2}{[S#2]}}
\newcommand{\cataloguescope}[1]{\paragraph{Recorded target.}#1}
\newcommand{\R}{\mathbb R}
\newcommand{\C}{\mathbb C}
\newcommand{\N}{\mathbb N}
\newcommand{\Z}{\mathbb Z}
\newcommand{\Q}{\mathbb Q}
\newcommand{\F}{\mathbb F}
\newcommand{\eps}{\varepsilon}
\begin{document}
% UnsolvedMath ID 30004281; source catalogue code OWR-17292-005
\setcounter{section}{30004280}
\section{Small-Weight Quantum Operations on Lattice Conformal Nets}\label{30004281}
{\small UnsolvedMath ID 30004281 · Mathematical Physics}
\subsection{Definitions and mathematical statement}

\paragraph{Shared notation.}$\mathbb N=\{1,2,\ldots\}$; the standard number fields have their usual meanings. All additional conventions are stated locally.


\paragraph{Definitions and conventions.}
For the source vacuum state $\omega$, a quantum operation is a compatible family $\phi_I$ of extreme normal unital completely positive, $\omega$-preserving maps on $\mathcal A(I)$, preserving local conformal symmetry and admitting a Markov adjoint in the same class. Composition is map composition. A nontrivial two-element hypergroup has an element $\phi\ne\operatorname{id}$ satisfying the displayed quadratic law, and $w$ is its hypergroup weight, not an arbitrary scalar attached to the identity map.

\paragraph{Target retained from the canonical source catalogue.}
\paragraph{Catalogue wording retained for traceability.}
For a lattice conformal net or a strongly local vertex-operator-algebra net $\mathcal A$, find $\phi\in\operatorname{QuOp}(\mathcal A)$ satisfying $$\phi\circ\phi=w^{-1}\operatorname{id}+(1-w^{-1})\phi$$ for arbitrarily small real $w>1$.

\paragraph{Source scope and qualification.}
Primary Problem 9 (PDF page 30, printed page 3082) asks for small positive real weights $w>1$. The catalogue adds arbitrarily small, which can be read as weights approaching $1$; that stronger quantifier was not verified in the primary source. The original small-weight construction question is preserved as source context, and the inherited stronger wording is explicitly not endorsed as a recovered theorem target. The nontriviality and actual weight convention exclude the vacuous identity operation.
\subsection{Short English statement}
Find conformal-net quantum operations with the stated quadratic composition law and weights arbitrarily close to one.
\subsection{Sources}
\begin{itemize}
\item[\hypertarget{source:30004281:1}{S1}] ULAM.AI and the UnsolvedMath Contributors, UnsolvedMath, record 30004281 (OWR-17292-005). Catalogue curation under CC BY 4.0; definitions and English exposition added locally.
\url{https://huggingface.co/datasets/ulamai/UnsolvedMath}.
\item[\hypertarget{source:30004281:2}{S2}] Subfactors and Applications (2019), Oberwolfach Reports, DOI 10.4171/owr/2019/49, pp. 3053--3140. Primary report accessed 2026-10-11; relevant PDF pages 29, 30.
\url{https://doi.org/10.4171/owr/2019/49}.
\item[\hypertarget{source:30004281:3}{S3}] Bischoff, generalized orbifold construction for conformal nets (catalogue-listed context; not independently consulted in this preparation).
\url{https://arxiv.org/abs/1608.00253}.
\item[\hypertarget{source:30004281:4}{S4}] Osborne et al., numerical evidence for a Haagerup CFT (catalogue-listed context; not independently consulted in this preparation).
\url{https://authors.library.caltech.edu/records/mgvpx-37124}.
\end{itemize}
\label{end:30004281}
\end{document}
